\chapter{Kết luận}

\section{Kết luận}
Luận văn này đã tập trung nghiên cứu và triển khai khung lý thuyết Đồ thị các
Tập lồi (Graphs of Convex Sets - GCS) để giải quyết bài toán hoạch định chuyển
động cho hệ thống robot trong các môi trường có cấu trúc hình học phức tạp. Về
mặt lý thuyết, chúng tôi đã xây dựng một quy trình hoàn chỉnh từ giai đoạn phân
hoạch không gian tự do thành hợp của các tập lồi, đến việc tham số hóa quỹ đạo
bằng đường cong Bézier nhằm đảm bảo tính liên tục của các trạng thái động học
cấp cao như vận tốc và gia tốc. Điểm cốt lõi của phương pháp nằm ở việc vận
dụng kỹ thuật nới lỏng phối cảnh để chuyển đổi bài toán tối ưu hóa lồi nguyên
hỗn hợp (MICP) vốn có độ phức tạp tính toán rất lớn thành bài toán Lập trình
nón bậc hai (SOCP) có thể giải quyết hiệu quả trong thời gian đa thức.

Thông qua các kết quả thực nghiệm định lượng trên môi trường 2D và hệ thống mê
cung đa tầng, luận văn đã chứng minh được tính khả thi và hiệu quả vận hành của
khung làm việc GCS. Trong các kịch bản môi trường 2D, việc phân tích đối chiếu
giữa các kỹ thuật phân hoạch không gian cho thấy thuật toán Phân hoạch lồi xấp
xỉ (ACD) mang lại sự cân bằng tối ưu giữa số lượng vùng lồi và thời gian hội tụ
của bộ giải so với phương pháp Lớp phủ Clique Khả kiến (VCC). Đặc biệt, kết quả
tính toán trên môi trường mê cung có độ phức tạp tổ hợp cao đã làm nổi bật sự
đánh đổi giữa hai mô hình: Linear GCS cho phép phản hồi cực nhanh về mặt hình
học, trong khi Bezier GCS mặc dù yêu cầu chi phí tính toán cao hơn nhưng lại
tạo ra các quỹ đạo mượt mà, đảm bảo các ràng buộc vi phân liên tục $C^2$ và khả
thi về mặt vật lý cho robot.

Một đóng góp quan trọng của nghiên cứu này là việc đề xuất và triển khai thành công mô hình bài toán Điều khiển Tối ưu (OCP) kết hợp Multiple Shooting và hàm cản log-barrier trên nền tảng của GCS. Thay vì chỉ tạo ra quỹ đạo hình học như GCS nguyên bản, hướng tiếp cận đề xuất đã giải quyết triệt để vấn đề vi phạm biên (hiện tượng đi xuyên vùng) thông qua áp dụng hàm phạt logarit. Đồng thời, kỹ thuật Multiple Shooting cho phép chia nhỏ và tối ưu hóa động lực học hệ thống một cách trơn tru, đảm bảo quỹ đạo thu được vừa tối ưu về mặt hình học, vừa khả thi tuyệt đối về mặt vật lý. Mặc dù trong phạm vi của báo cáo này, chúng tôi tập trung chủ yếu vào hệ thống robot trong môi trường tĩnh, nhưng các dữ liệu thực nghiệm định lượng thu được đã minh chứng tính hiệu quả vượt trội của phương pháp đề xuất.

\section{Hướng nghiên cứu trong tương lai}

Dựa trên những kết quả thực nghiệm đạt được, nghiên cứu đề xuất các hướng phát triển tiếp theo nhằm tối ưu hóa hiệu năng và mở rộng khả năng ứng dụng của đồ án như sau:

\begin{itemize}
    \item \textbf{Tối ưu hóa hiệu suất tính toán thời gian thực (Real-time MPC):} Nghiên cứu sẽ tập trung vào việc áp dụng phương pháp của đồ án trong mô hình Điều khiển Dự báo Mô hình (Model Predictive Control - MPC). Việc khai thác khả năng warm-start của Multiple Shooting sẽ giúp hệ thống đáp ứng tốt hơn các yêu cầu hoạch định quỹ đạo trực tuyến (online planning) cho robot.

    \item \textbf{Mở rộng cho môi trường động (Dynamic Environments):} Khung làm việc hiện tại chủ yếu giải quyết bài toán tĩnh. Hướng đi tiếp theo là mở rộng đồ thị tập lồi theo cả trục thời gian (Spatiotemporal GCS) để nhận diện và lẩn tránh các vật cản di động.

    \item \textbf{Nâng cao thuật toán phân hoạch không gian an toàn:} Để xử lý các môi trường có vật thể phi tuyến tính phức tạp trong không gian đa chiều, nghiên cứu dự kiến sẽ tích hợp các kỹ thuật phân hoạch trực tiếp trong không gian cấu hình ($C$-space) như C-IRIS. Điều này cho phép mở rộng áp dụng hệ thống cho các cánh tay robot công nghiệp có số bậc tự do (DOF) lớn.

    \item \textbf{Cải tiến hàm phạt Log-barrier:} Đánh giá thêm các dạng hàm phạt hoặc các phương pháp interior-point nâng cao nhằm giảm số lần lặp (iterations) cần thiết khi điều chỉnh trọng số $\mu$, qua đó tối ưu thêm chi phí tính toán của bài toán NLP.
\end{itemize}

\newpage
Để cụ thể hóa các hướng nghiên cứu trên, bảng \ref{tab:graduation-plan} trình bày kế hoạch chi tiết trong 15 tuần tới nhằm hoàn thiện đồ án tốt nghiệp, tập trung vào việc kiểm thử và đánh giá các cải tiến đề xuất.

\begin{table}[htbp]
    \centering
    \caption{Kế hoạch chi tiết Đồ án Tốt nghiệp (15 tuần)}
    \label{tab:graduation-plan}
    \footnotesize
    \setlength{\tabcolsep}{4pt}
    \renewcommand{\arraystretch}{1.3} % Tăng độ giãn dòng để các ô đa dòng dễ đọc hơn
    \begin{tabular}{|>{\raggedright\arraybackslash}m{2.2cm}|
        >{\raggedright\arraybackslash}m{6cm}|
        >{\centering\arraybackslash}m{1cm}|
        >{\centering\arraybackslash}m{1cm}|
        >{\centering\arraybackslash}m{1cm}|
        >{\centering\arraybackslash}m{1cm}|
        >{\centering\arraybackslash}m{1cm}|}
        \hline
        \textbf{Giai đoạn} & \textbf{Nhiệm vụ chi tiết}                                                          & \textbf{W1--3} & \textbf{W4--6} & \textbf{W7--9} & \textbf{W10--12} & \textbf{W13--15} \\ \hline

        % Giai đoạn 1
        \multirow{4}{2.4cm}{\textbf{Tối ưu thuật toán}}
                           & (1) Kiểm thử và tích hợp kỹ thuật Multiple Shooting vào GCS Framework               & X              &                &                &                  &                  \\ \cline{2-7}
                           & (2) Triển khai phân hoạch $C$-space bằng C-IRIS/IRIS-NP để xử lý vật cản phi tuyến  & X              & X              &                &                  &                  \\ \hline

        % Giai đoạn 2
        \multirow{4}{2.4cm}{\textbf{Hiện thực mô hình}}
                           & (3) Thiết kế hệ thống điều khiển robot tích hợp khung BezierGCS                     &                & X              & X              &                  &                  \\ \cline{2-7}
                           & (4) Xây dựng thuật toán nhận diện và ánh xạ mê cung thực tế sang đồ thị các tập lồi &                &                & X              & X                &                  \\ \hline

        % Giai đoạn 3
        \multirow{4}{2.4cm}{\textbf{Thực nghiệm}}
                           & (5) Thử nghiệm robot giải mê cung, đo lường thời gian giải và sai số quỹ đạo        &                &                &                & X                & X                \\ \cline{2-7}
                           & (6) Đánh giá đối chứng hiệu năng với các thuật toán lấy mẫu RRT* và PRM             &                &                &                &                  & X                \\ \hline

        % Giai đoạn 4
        \multirow{2}{2.4cm}{\textbf{Hoàn thiện}}
                           & (7) Phân tích dữ liệu, đánh giá tính hội tụ của nới lỏng phối cảnh và viết luận văn &                &                &                &                  & X                \\ \hline
    \end{tabular}
\end{table}